\documentclass[b5paper,10pt,twoside]{book}
\usepackage[T1]{fontenc}
\usepackage[utf8]{inputenc}
\usepackage[spanish]{babel}
\usepackage[super]{nth}
\usepackage{
    amsmath,
    amsfonts,
    amssymb,
    amsthm,
    imakeidx,
    graphicx,
    polynom,
    % subcaption,
    lipsum,
    lmodern,
    systeme,
    float,
    ifthen,
    mathrsfs,
    units,
    % enumerate,
    enumitem,
    multicol,
    geometry,
    mdframed,
    setspace,
    subcaption,
    anyfontsize,
    mathtools,
    fancyhdr,
    booktabs,
    multirow,
    csquotes,
    microtype,
    bookmark,
    emptypage,
    blkarray,
    aliascnt,
    dsfont,
    comment,
    tabularx,
    % mathds,
    titlesec,
    cancel,
    currfile,
    xcolor,
}
\usepackage[perpage]{footmisc}

\setenumerate{topsep=0pt}

\setlength{\parskip}{0.5\baselineskip}%

% \fontseries{b}\selectfont
% \renewcommand{\bfdefault}{b}

\titleformat{\chapter}
  {\normalfont\Huge}{\Huge\thechapter}{15pt}{\Huge}

% \titleformat{\section}
%   {\normalfont\large}{\large\thechapter}{15pt}{\large}

% \titleformat{\section}
%   {\normalfont\Huge}{\Huge\thechapter}{15pt}{\Huge}

% \onehalfspacing
\setstretch{1.05}
\geometry{
    margin=1in,
    % headheight=50pt,
    % headsep=12pt,
    % showframe
    % textwidth=5in,
}
% \decimalpoint
\raggedbottom
\pagestyle{plain}
% \fancyhf{}
% % \fancyheadoffset{1.0mm}
% \lhead{\includegraphics[width=0.35\textwidth]{logoECMC.png}}
% \rhead{\includegraphics[width=0.25\textwidth]{logoYachay.png}}
\fancyfoot[RO,LE]{\thepage}

\usepackage[delims=\lbrack\rbrack]{spalign}
\let\matrix=\spalignmat
\let\Amatrix=\spalignaugmat
\let\AAmatrix=\spalignaugmatn
\newcommand{\Dmatrix}[1]{\spaligndelims\vert\vert\spalignmat{#1}}

% \fancyhead[RE]{\scshape{\nouppercase{\leftmark}}}
% \fancyhead[LO]{\scshape{\nouppercase{\rightmark}}}
\renewcommand{\headrulewidth}{0pt}
\graphicspath{{images/}}
\bookmarksetup{numbered}
\makeatletter
 \renewcommand\Hy@numberline[1]{#1. }
\makeatother
\usepackage{hyperref}
\hypersetup{%
    colorlinks = {true},
            linkcolor = {blue},
            urlcolor  = {blue},
            citecolor = {blue},
            anchorcolor = {blue},
    pdftitle    ={},
    pdfauthor   ={K. Christian Chávez},
    pdfsubject  ={Abstract Algebra},
    pdfcreator  ={LaTeX},
    pdfproducer ={},
    % pdfkeywords={}{}{},
    % hidelinks,
    pdfpagemode={UseNone},
}
\newcolumntype{Y}{>{\centering\arraybackslash}X}
\renewcommand\tabularxcolumn[1]{m{#1}}
% \renewcommand\arraystretch{1.5}

\newcommand{\p}{\cdot}
\newenvironment{solution}{\noindent\textbf{\textsl{Solution}.}}{\hfill$\square$}
% \newenvironment{remark}{\noindent\textbf{\textsl{Remark.}}}{}
% \theoremstyle{definition}
% \newtheorem{problem}{Problem}%[chapter]

% \setlength\parindent{0pt}

% Nuevos comandos
\let\originalleft\left
\let\originalright\right
% \renewcommand{\left}{\mathopen{}\mathclose\bgroup\originalleft}
% \renewcommand{\right}{\aftergroup\egroup\originalright}
\newcommand{\N}{\mathds{N}}
\newcommand{\Z}{\mathds{Z}}
\newcommand{\R}{\mathds{R}}
\newcommand{\Q}{\mathds{Q}}
\newcommand{\II}{\mathds{I}}
\newcommand{\K}{\mathds{K}}
\newcommand{\PP}{\mathds{P}}
\newcommand{\C}{\mathds{C}}
\newcommand{\T}{\mathscr{T}}
\newcommand{\I}{\mathcal{I}}
\newcommand{\E}{\mathscr{E}}
\newcommand{\B}{\mathcal{B}}
\newcommand{\F}{\mathscr{F}}
\newcommand{\U}{\mathscr{U}}
\newcommand{\LL}{\mathscr{L}}
\newcommand{\OO}{\mathscr{O}}
\newcommand{\Fam}{\mathscr{F}}
\newcommand{\Nei}{\mathscr{N}}
\newcommand{\Pt}{\mathscr{P}}
% \newcommand{\ES}{\text{\upshape\O}}
% \renewcommand{\ES}{\cancel{\mathrm{O}}}
\newcommand{\ES}{\varnothing}
\newcommand{\TopS}{\left( X, \T \right)}
\newcommand{\bcap}{\bigcap}
\newcommand{\bcup}{\bigcup}
\newcommand{\dif}{\backslash}
\newcommand{\Solution}{\noindent\textbf{Solution.}}
% \let\oldbigcup\bigcup
% \renewcommand{\bigcup}{\boldsymbol{\oldbigcup}} 

\newcommand{\ev}{\textsc{ev}}
\newcommand{\ii}{\hat{\imath}}
\newcommand{\jj}{\hat{\jmath}}
\newcommand{\kk}{\hat{k}}
\newcommand{\vv}{\mathbf{v}}
\newcommand{\va}{\mathbf{a}}
\newcommand{\vb}{\mathbf{b}}
\newcommand{\vd}{\mathbf{d}}
\newcommand{\vx}{\mathbf{x}}
\newcommand{\vy}{\mathbf{y}}
\newcommand{\vu}{\mathbf{u}}
\newcommand{\vzr}{\mathbf{0}}
\newcommand{\pp}{\boldsymbol{\cdot}}
\newcommand{\AND}{\quad\text{and}\quad}
\newcommand{\seq}{\subseteq}
\newcommand{\sep}{\supseteq}
\newcommand{\nseq}{\nsubseteq}
\newcommand{\then}{\implies}
\newcommand{\abs}[1]{\left|#1\right|}
\newcommand{\Abs}[1]{\left|\!\left|#1\right|\!\right|}
\newcommand{\prt}[1]{\left(#1\right)}
\newcommand{\End}{\hfill\(\square\)}
\newcommand{\Mod}[1]{\ \left(\mathrm{mod}\ #1\right)}
\newcommand{\floor}[1]{\left\lfloor #1 \right\rfloor}
\newcommand{\ceil}[1]{\left\lceil #1 \right\rceil}
\renewcommand\labelitemi{$\bullet$}
\newcommand{\interior}[1]{#1^\circ}
\renewcommand\labelitemi{$\bullet$}
\renewcommand{\labelitemii}{$ \circ $}
\renewcommand{\mathbb}[1]{\mathds{#1}}
\renewcommand*{\thefootnote}{\fnsymbol{footnote}}

\DeclareMathOperator{\lcm}{lcm}
\DeclareMathOperator{\adj}{adj}
\DeclareMathOperator{\sen}{sen}
\DeclareMathOperator{\proy}{proy}
\DeclareMathOperator{\tr}{tr}
\DeclareMathOperator{\spn}{span}
\DeclareMathOperator{\im}{Im}
\DeclareMathOperator{\dd}{d }
\DeclareMathOperator{\card}{card}
\DeclareMathOperator{\ddd}{d\!}
% \DeclareMathOperator{\y}{y}
\DeclareMathOperator{\rank}{rank}
\DeclareMathOperator{\nul}{nul}
\DeclareMathOperator{\Int}{int}
\DeclareMathOperator{\cla}{cla}

% Añadidos después
\newcommand{\vect}[1]{\left( {#1}_{1},\dots,{#1}_{n} \right)}

\usepackage{amsthm}
\newtheoremstyle{uptheorem} 
{0.75cm}{0.75cm}{\upshape}{}{\bfseries}{.}{ }{} 
\newtheoremstyle{sltheorem} 
{0.75cm}{0.3cm}{\slshape}{}{\bfseries}{.}{ }{} 
% \theoremstyle{sltheorem}
\theoremstyle{uptheorem}
\newtheorem{theorem}{Problema}%[chapter]
\newtheorem{definition}{Definition}%[chapter]
\newcommand{\definitionautorefname}{Definition}
\newtheorem{lemma}{Lemma}%[theorem]
\newcommand{\lemmaautorefname}{Lemma}
\newtheorem{corollary}{Corollary}%[theorem]
\newcommand{\corollaryautorefname}{Corollary}
\newtheorem{corolario}{Corolario}[theorem]
\newcommand{\corolarioautorefname}{Corolario}
\theoremstyle{uptheorem} 
\newtheorem{example}{Example}%[chapter]
  \newaliascnt{examples}{example}
  \newtheorem{examples}[examples]{Examples}
  \aliascntresetthe{examples}
  \def\exampleautorefname{Example}
  \def\examplesautorefname{Examples}

\newcommand{\notaautorefname}{Nota}

\newtheoremstyle{uptheoremv2}{0.7cm}{}{\upshape}{}{\bfseries}{.}{ }{} 
\theoremstyle{uptheoremv2} 
\newtheorem{problem}[theorem]{Problema}
\newcommand{\problemautorefname}{Problem}
\newtheoremstyle{uptheoremv3}{}{}{\upshape\small}{}{\bfseries}{.}{ }{} 
\theoremstyle{uptheoremv3}
\newtheorem{remark}{Remark}[chapter]

\numberwithin{figure}{section}
\numberwithin{table}{section}
% \numberwithin{equation}{chapter}

% \usepackage[backend=biber]{biblatex}
% \bibliography{bib.bib}
% \nocite{*}

\color{black!85}

\begin{document}
\def\thetitle{Tarea 01 -- Parte 2}


% \pagebreak
% \vspace*{-1.5cm}
% \begin{center}
%     \includegraphics[height=0.075\textheight]{images/LogoClubMate.pdf} 
%         % \hspace{1cm}
%     %     \hfill
%     % \includegraphics[height=0.075\textheight]{images/LogoECMC.pdf}
% \end{center}

% \phantomsection\pdfbookmark{\currfilebase}{\currfilebase}

% \begin{center}
%     {\LARGE
%     Topología \\ Curso Vacacional Enero 2025\\
%     % Teaching Assistant
%     \vspace{0.5cm}
%     \textbf{\thetitle{}}}

%     % \emph{I hear, I forget;
%     % I see, I remember;
%     % I do, I understand.}
    
%     Christian Chávez\\
%     \today
% \end{center}

% % \setcounter{problem}{0}


%=================================

\begin{center}
\begin{minipage}{0.45\textwidth}
    \includegraphics[height=0.075\textheight]{images/LogoClubMate.pdf}
\end{minipage}%
\hfill
\begin{minipage}{0.5\textwidth}
    \begin{flushright}
        {\LARGE Topología} \\ Curso Vacacional -- Enero 2025 \\
        Instructor: Christian Chávez \\
    \end{flushright}
\end{minipage}

\vspace{\baselineskip}
{\huge\thetitle} \\
% \vspace{0.5cm}

\today\\
Fecha de entrega: \fechaentrega{}

\end{center}



% \section*{1.\enspace Espacios topológicos}


% \begin{problem}
% Sea $X$ un conjunto infinito y $p \in X$. Muestra que
% \begin{align*}
% \mathcal{T}_1 &= \{ A \subseteq X \mid A = \varnothing \, \text{ o } \, X \setminus A \text{ es finito} \}\quad\text{y}\\
% \mathcal{T}_2 &= \{ A \subseteq X \mid A = \varnothing \, \text{ o } \, p \in A \}
% \end{align*}
% son topologías sobre $X$.

% \end{problem}



% \begin{problem}\hfill\null
% \begin{enumerate}[label=(\roman*)]
% \item Sea $X = \{a, b, c, d\}$ dotado de la topología
% \[
% \mathcal{T} = \{\varnothing, X, \{a\}, \{b\}, \{a, c\}, \{a, b\}, \{a, b, c\}\}.
% \]
% Encuentra una base y una subbase para $\mathcal{T}$.

% \item ¿Por qué una topología es una base para ella misma?

% \item Muestra que \(\mathcal{B}  = \{(a, b) \subseteq \mathbb{R} \mid a, b \in \mathbb{Q}\}\)  
% es una  base en $\mathbb{R}$ y que además $\mathcal{T}_{\mathcal{B}} = \mathcal{T}_{\text{usual}}$. (Ayuda: usa el hecho de que todo real es el límite de una secuencia de racionales.)
% \end{enumerate}


% \end{problem}
 

% % \begin{problem}
% %     % Muestra que la topología generada por una base es única.
% % \end{problem}



% \section*{2.\enspace Espacios métricos}

% \begin{problem}
% \begin{enumerate}[label=(\roman*)]
% \item Muestra que la suma/el máximo de métricas es una métrica. En otras palabras, si \(d\) y \(\rho\) son métricas sobre un conjunto \(M\), entonces \(d+\rho\) definida por \((d+\rho)(x,y) = d(x,y) + \rho(x,y)\) es una métrica sobre \(M\). Similarmente, \(\xi(x,y) = \max\left\{ d(x,y) , \rho(x,y) \right\}\) es una métrica sobre \(M\). 

% \item Muestra que las siguientes funciones son métricas en $\mathbb{R}^n$.
% \begin{enumerate}[label=(\alph*)]
%     \item $d(x, y) = \max\{|x_i - y_i| : 1 \leq i \leq n\}$.
%     \item $\rho(x, y) = |x_1 - y_1| + \cdots + |x_n - y_n|$.
% \end{enumerate}
% \end{enumerate}
% \end{problem}


% \begin{problem}
% Sea $d$ una métrica en un conjunto $X$.

% \begin{enumerate}[label=(\roman*)]
% \item Muestra que la función ${D} : X \times X \to [0, +\infty)$ definida por 
% \[
% (x, y) \mapsto \frac{d(x, y)}{1 + d(x, y)}
% \]
% es una métrica en $X$.

% \item Muestra que $d$ y ${D}$ son equivalentes.
 
% \item Muestra que 
% \[
% \rho : X \times X \to [0, +\infty) \; : \;  (x, y) \mapsto \min\{1, d(x, y)\}
% \]
% es una métrica en $X$.

% \item Supón que $f : [0, +\infty) \to [0, +\infty)$ es una función monótona creciente tal que $f(0) = 0$ y 
% \[
% f(x + y) \leq f(x) + f(y)
% \]
% para todo \(x,y\geq 0\).
% Muestra que $f \circ d$ es una métrica en $X$.
% \end{enumerate}
    
% \end{problem}

\section*{3.\enspace Posición de un punto respecto  a un conjunto}


\begin{problem}
    Sea \(A\)  un subconjunto  de un espacio topológico \(X\).
    Verifica que \(A^\circ \subset A\) y \(A\subset \overline{A}\).
    Muestra que \((A^\circ)^\circ = A^\circ\) y \(\overline{\overline{A}} = \overline{A}\). 

\end{problem}


\begin{problem}
    Sean \(A\) y \(B\) un subconjuntos de un espacio topológico.
    Demuestra que 
    \begin{enumerate}[label=(\roman*)]
        \item si \(A\subset B\), entonces \(A^\circ \subset B^\circ\),
        \item si \(A\subset B\), entonces \(\text{Ext}\, {B} \subset \text{Ext}\, {A}\),
        \item si \(A\subset B\), entonces \(\overline{A} \subset \overline{B}\),
        \item \((A\cap B)^\circ  = A^\circ \cap B^\circ \) y \(\overline{A\cup B}=\overline{A}\cup \overline{B}\).
    \end{enumerate}
\end{problem}
    

\begin{problem}
Sea \(A\) un subconjunto de un espacio topológico.
Demuestra que 
\begin{enumerate}[label=(\roman*)]
    \item \(A\) es abierto si y solo si \(A = A^\circ\),
    \item \(A\) es cerrado si y solo si \(A = \overline{A}\),
    \item \(A\) es cerrado si y solo si \(\partial A \subset  {A}\),
    \item \(A\) es cerrado si y solo si \( A' \subset  {A}\).
\end{enumerate}
\end{problem}


\begin{problem} 
    Sean \(A\) y \(B\) un subconjuntos de un espacio topológico.
\begin{enumerate}[label=(\roman*)]
    \item Demuestra que \(\partial A = \partial (X\setminus A)\).
    \item Demuestra que \(\partial A = \overline{A}\cap \overline{X\setminus A}   \).
    \item Encuentra contraejemplos para \((A\cup  B)^\circ  = A^\circ \cup  B^\circ \) y \(\overline{A\cap B}=\overline{A}\cap \overline{B}\).
    \item Prueba que si \(A\) es un subespacio de \(X\) y \(B\subset A\), entonces \(\text{Cl}_A \, B = (\text{Cl}_X B)\cap A\).
    Aquí, \(\text{Cl}_A\) denota la clausura respecto a \(A\), tomado como espacio topológico.
\end{enumerate}
\end{problem}




\end{document}

